\chapter{Rule Properties and Modification Logic}
\label{ch:rule-properties}

The \fn{declare} block of a rule sets properties that control when and
how the rule fires. \Morendo\ accepts these, in any order:

\begin{center}
\begin{tabularx}{\textwidth}{llL}
\toprule
Property & Value & Effect \\
\midrule
\fn{salience} & integer & priority on the agenda; default 100, higher fires first \\
\fn{auto-focus} & boolean & bring the rule's module into focus when it activates \\
\fn{rule-version} & symbol or number & a label shown by \fn{(rules)} and \fn{ppdefrule} \\
\fn{remember-alpha} & boolean & keep alpha memories for the rule (default true) \\
\fn{hashed-memory} & boolean & hash the rule's join memories (default true) \\
\fn{effective-date} & string & the rule activates only after this time \\
\fn{expiration-date} & string & and only before this time \\
\fn{chaining-direction} & \fn{forward} or \fn{backward} & recorded on the rule; the engine chains forward \\
\fn{no-agenda} & boolean & fire immediately, bypassing the agenda \\
\fn{temporal-activation} & boolean & see Chapter~\ref{ch:temporal} \\
\bottomrule
\end{tabularx}
\end{center}

\fn{ppdefrule} prints the declaration with every property that differs
from a bare rule, in the form the parser accepts, so it is the quickest
way to see the defaults.

\section{Salience}

Salience is an integer literal, possibly negative. Activations of higher
salience fire first regardless of strategy. Section~\ref{sec:tutorial-salience}
shows the effect; a common pattern is a low-salience ``cleanup'' rule
that retracts a control fact after every other rule for it has fired.

\section{Event-driven rules: no-agenda}
\label{sec:no-agenda}

\begin{lstlisting}
(defrule make_path (declare (no-agenda TRUE))
   ?context <- (context (name "make_path"))
   (seating (id ?id) (pid ?sPID) (path_done FALSE))
   (path (id ?sPID) (name ?n1) (seat ?s))
   (not (path (id ?id) (name ?n1)))
=>
   (assert (path (id ?id) (name ?n1) (seat ?s))))
\end{lstlisting}

A \fn{no-agenda} rule fires the moment a complete match reaches its
terminal node, inside the \fn{assert}, \fn{modify} or \fn{retract} that
completed the match, without waiting for \fn{fire}. Salience and
strategy do not apply to it. This turns the engine into an event
processor: as facts stream in, the reactions run at once, while ordinary
rules still accumulate on the agenda for the next \fn{fire}. The Manners
benchmark uses it for the rule above, which must propagate before any
other rule looks at the seating, and the engine uses it for the rule that
feeds each MOLAP cube (Chapter~\ref{ch:molap}).

Because such a rule runs inside another operation, its actions should be
short and must not assume any particular order relative to the agenda.

\section{Effective and expiration dates}
\label{sec:rule-dates}

\begin{lstlisting}
(defrule live
  (declare (effective-date "09/01/2026 00:00") (expiration-date "12/31/2026 23:59"))
  (order (total ?t))
=>
  (printout t "live fired" crlf))
(defrule dead
  (declare (effective-date "01/01/2010 00:00") (expiration-date "03/31/2010 23:59"))
  (order (total ?t))
=>
  (printout t "dead fired" crlf))
\end{lstlisting}

Asserting an \fn{order} and firing prints only \fn{live fired}. A rule
with an expiration date is given a terminal node that checks the clock: a
match is put on the agenda only if the current time is after the
effective date and before the expiration date. Matches outside the window
are dropped, and they are not reconsidered when the window opens later;
the fact has to be asserted or modified again. Retractions are always
honoured. A rule with only an effective date and no expiration date is
not time-checked at all.

Dates are written \fn{MM/dd/yyyy HH:mm} (month, day, year, hour and
minute, 24-hour clock) and interpreted in the JVM's default time zone.

\section{Remembering alpha memories}

By default every alpha node keeps the facts that passed it, so that a
join can be recomputed when the other side changes. \fn{(remember-alpha
FALSE)} tells the compiler to drop that memory for the rule's nodes,
trading memory for repeated evaluation. It is an optimization for rules
whose patterns are cheap to re-test and match many facts, and it only
affects nodes not shared with other rules.

\section{Auto-focus and modules}
\label{sec:auto-focus}

A rule declared \fn{(auto-focus TRUE)} brings its module into focus the
moment one of its activations reaches the agenda, remembering the
previous focus. \fn{fire} drains the module in focus; when that module's
agenda is empty it returns to the module that was in focus before and
continues there, so a rule set can hand control to a specialised module
and get it back without an explicit \fn{set-focus}:

\begin{lstlisting}[style=shell]
Morendo> (defmodule ACCT)
true
Morendo> (deftemplate ACCT::u (slot a))
true
Morendo> (defrule ACCT::r (u (a ?x)) => (printout t "module rule " ?x crlf))
true
Morendo> (defrule ACCT::auto (declare (auto-focus TRUE)) (u (a ?x)) => (printout t "auto-focus rule " ?x crlf))
true
Morendo> (set-focus MAIN)
ACCT
Morendo> (assert (ACCT::u (a 2)))
2
Morendo> (focus)
ACCT
Morendo> (fire)
module rule 2
auto-focus rule 2
2
Morendo> (focus)
MAIN
\end{lstlisting}

\section{Other properties}

\fn{rule-version} and \fn{chaining-direction} are recorded on the rule
and shown by \fn{(rules)} and \fn{ppdefrule}; the engine itself always
chains forward, so \fn{chaining-direction} is documentation for tools
built on the Java API. \fn{hashed-memory} controls whether the rule's
join memories are hashed by their bound values (the default) or scanned.

\section{Modification logic}
\label{sec:modification-logic}

A rule's actions run when a match appears. Nothing runs when a match
\emph{disappears}: if a fact the rule used is retracted or modified, the
activation is simply withdrawn (or, if it already fired, nothing happens
at all). For rules that maintain aggregates that is a problem: a rule
that adds a transaction to a running total has no way to subtract it when
the transaction is cancelled, short of a second rule and a status slot to
keep the two from fighting.

\Morendo\ extends \fn{defrule} with a second action block, introduced by
\fn{<<=}, that runs when a fact of a match that has already fired is
modified or retracted:

\begin{lstlisting}
(deftemplate item (slot name))
(deftemplate seen (slot name))
(defrule track
  ?i <- (item (name ?n))
=>
  (printout t "add " ?n crlf)
  (assert (seen (name ?n)))
<<=
  (printout t "remove " ?n crlf))
\end{lstlisting}

\begin{lstlisting}[style=shell]
Morendo> (assert (item (name "a")))
2
Morendo> (fire)
add a
1
Morendo> (retract 2)
true
Morendo> (fire)
remove a
1
Morendo> (assert (item (name "b")))
4
Morendo> (retract 4)
true
Morendo> (fire)
0
\end{lstlisting}

The \fn{=>} actions run when the match appears; the \fn{<<=} actions are
the compensating step, run with the same variable bindings when a fact of
the match is retracted after the rule has fired. A rule with modification
actions is compiled with a special terminal node: on retraction, if the
activation is still waiting on the agenda it is withdrawn as usual
(item \fn{b} above), and if the rule has already fired a
\emph{modification activation} is placed on the agenda instead, which runs
the \fn{<<=} block at the next \fn{fire}. The design, with the
aggregate-maintenance rule that motivated it, is written up in
\file{doc/modification\_logic.pdf}.

Modification actions fire once per retracted match, with the bindings the
original activation had, including the retracted fact's values. They are
the mechanism behind the cube-maintenance rules the engine generates.
